\documentclass[11pt,a4paper]{article}
\usepackage[T1]{fontenc}
\usepackage{lmodern}
\usepackage[margin=24mm,headheight=15pt]{geometry}
\usepackage{amsmath,amssymb,amsthm,mathtools,booktabs,array,tabularx}
\usepackage{microtype,fancyhdr,xcolor}
\usepackage[unicode]{hyperref}
\definecolor{ink}{RGB}{28,61,79}
\hypersetup{colorlinks=true,linkcolor=ink,urlcolor=ink,citecolor=ink,
 pdftitle={SRG Amplitude Equality: Conditional Persistence and Minimal Splitting in Tower II},
 pdfauthor={Hilmir Frímann Halldórsson}}
\pagestyle{fancy}\fancyhf{}
\fancyhead[L]{\small SRG Amplitude Equality}
\fancyhead[R]{\small v0.1 | Manuscript for review}
\fancyfoot[C]{\thepage}
\setlength{\parindent}{0pt}
\setlength{\parskip}{5pt plus 1pt minus 1pt}
\setlength{\emergencystretch}{2em}
\clubpenalty=10000 \widowpenalty=10000
\allowdisplaybreaks[1]
\newcolumntype{Y}{>{\raggedright\arraybackslash}X}
\theoremstyle{definition}
\newtheorem{theorem}{Theorem}[section]
\newtheorem{proposition}[theorem]{Proposition}
\theoremstyle{remark}\newtheorem{remark}[theorem]{Remark}
\DeclareMathOperator{\sgn}{sgn}
\newcommand{\R}{\mathbb R}
\newcommand{\ind}{\mathbf 1}
\numberwithin{equation}{section}
\title{\textbf{SRG Amplitude Equality}\\[4pt]
\large Conditional Persistence and Minimal Splitting in Tower II}
\author{Hilmir Frímann Halldórsson}
\date{Version v0.1 --- 12 October 2026\\[3pt]\small Manuscript and proposed publication payload awaiting review}
\begin{document}
\maketitle
\begin{abstract}
When equal particle amplitudes enter the historical Tower II composition, do they leave equal?
The answer depends on the processing branches and frequencies, as well as on compression.
We reproduce the accepted exact-real analysis of the unchanged source.
Equal frequencies and active rows preserve amplitude equality at each pass, even when
the normalized distance weights differ; persistence through several passes is conditional
on the hypotheses continuing to hold. Three collinear particles give a minimal
common-frequency geometric counterexample: a close pair receives the sine stage,
while an isolated particle skips it. A different minimal example uses the source's actual
two-band assignment: two active particles split because their sine factors differ,
although their complete enhancement increments agree. We give the branch proofs,
strict guard bounds and compressor qualifications, including signed and zero states.
These are amplitude results for specified incoming states, not trajectory, full-state
capture or physical-validation results.
\end{abstract}

\section{Three particles, one amplitude, two processing branches}\label{sec:opening}
Place three particles on a line, at coordinates $0$, $1/2$ and $3$.
Give every particle amplitude $1/5$, frequency $0.244$ and a common neighbor horizon
of radius $1$. Evaluate the next pass at its incremented clock $t=1$.
The source's horizon is strict: a neighbor must be closer than $1$.

\begin{center}
\begin{tabular}{@{}cccl@{}}
\toprule Particle & Position & Selected neighbors & Amplitude processing\\
\midrule
1 & $(0,0,0)$ & 2 & average, sine, enhancement\\
2 & $(1/2,0,0)$ & 1 & average, sine, enhancement\\
3 & $(3,0,0)$ & none & these three stages skipped\\
\bottomrule
\end{tabular}
\end{center}

The close pair averages two copies of the same number, so averaging itself changes
nothing. Both pair members then receive the same sine multiplier. The isolated
particle does not. A later common multiplier involving $H(t)$ applies to all three.
The compressor is on the same positive, unclipped branch for this example, so it
does not erase the difference. The pair ends approximately
$0.002058402961068474$ above the isolated particle; this decimal is a rounded
transcription of the corrected, accepted OC05 result, not a new evaluation.

This elementary picture separates two questions. Changing distances can change which
rows are processed; yet, while every row is processed and the frequencies agree,
changing normalized weights alone cannot separate equal amplitudes. A second example,
with two different source frequencies and no skipped row, separates amplitudes for
a different reason. We prove both statements below.

\section{The historical source contract}\label{sec:source}
We use the frozen Tower II text \cite{T}, not a replacement update.
The accepted analyses OC05 and OC07 \cite{OC05,OC07}, the earlier OC02 framework
\cite{G}, and the source-order addendum \cite{A} supply the mathematical provenance.
The line map in Appendix~\ref{app:sources} identifies the physically inspected source
definitions. Tower II is a historical model name; this paper is not the reserved
SRG-II/R0--R3 mechanism manuscript.

\subsection{Incoming data, selection and order}
Write $R_i\in\R$ for the incoming amplitude, $x_i\in\R^3$ for position and
$\omega_i>0$ for frequency. The source increments its clock before the stages
studied here. Thus $t$ below is that incremented clock, not the incoming clock; we restrict to $t\geq0$.
All statements assume finite source-valid data and well-defined finite intermediate
quantities. They concern the exact-real interpretation of the formulas, not an
assertion of bitwise equality in a floating-point execution.

The source selects up to five directed neighbors. With a horizon $\rho$, it explicitly
excludes self, retains $d_{ij}=\|x_i-x_j\|<\rho$, and takes up to the nearest five.
Without a horizon it uses the sorted-distance slice \texttt{[1:6]}.
For distinct positions this removes self; coincident ties must not be silently
identified with a universal explicit-self-exclusion rule.
Neighbor sets and distances are stored from the incoming positions. The spatial
motion stage occurs before the amplitude loop, but the amplitude loop uses those
stored sets and distances.

For selected neighbors $N_i$, define
\begin{equation}\label{eq:weights}
 u_{ij}=\frac{\exp[-c\log^2(\omega_i/\omega_j)]}{d_{ij}+10^{-8}},
 \qquad S_i=\sum_{j\in N_i}u_{ij},\qquad
 w_{ij}=\frac{u_{ij}}{S_i}.
\end{equation}
Here $c=3$ is the retained default. A row is skipped if $N_i$ is empty or
$S_i<10^{-10}$. Otherwise it is active; equality at the threshold is processed.
The fixtures below have $S_i>10^{-10}$ strictly, avoiding the boundary.

Let $\beta_T=0.15$ be the Tower amplitude mixing coefficient, distinct from any
position coefficient. Define
\begin{equation}\label{eq:gains}
 a_i(t)=1+0.01\sin(2\pi\omega_i t),\qquad
 b_t=1+0.05\,\frac{H(t)}{H_{\max}}.
\end{equation}
The precompression amplitude, after the common $H$ stage, is
\begin{equation}\label{eq:full}
 q_i=b_t
 \begin{cases}
 a_i(t)\left[(1-\beta_T)R_i+\beta_T\sum_{j\in N_i}w_{ij}R_j\right]
       +\displaystyle\sum_{j\in N_i}\eta_{ij}R_j,&i\text{ active},\\
 R_i,&i\text{ skipped}.
 \end{cases}
\end{equation}
The sine multiplies the mixed baseline \emph{before} enhancement is added.
The full $\eta_{ij}$, including every strict guard and tower-list term, is given
in Appendix~\ref{app:eta}. It is not a normalized distance weight.

Following the common $H$ multiplication, the source performs positional fusion
using the precompression amplitudes, with weights $|q_i|+10^{-8}$ and only
bands containing at least two particles. RPCO then compresses amplitudes.
Subsequent noise and velocity updates affect positions or velocities.
These positional stages do not change this same-pass amplitude formula;
they may affect the neighbor geometry of later passes.

\subsection{The retained H and its bound}\label{sec:H}
The source uses
\begin{equation}\label{eq:H}
 H(t)=2e^{-t/2}\cos t+e^{-3t/10}\sin t
               +\tfrac32 e^{-7t/10}+e^{-2t/5}.
\end{equation}
Its normalization samples $50{,}000$ points on $[0,100]$, including zero, takes
the largest absolute value, and adds $10^{-12}$. The following accepted exact-real
bound \cite{G,OC05} explains the mathematical normalization without rerunning that
sampling calculation.

For $t\geq0$, $|\sin t|\leq t$ gives
\[
 |H(t)|\leq B(t)+\tfrac32e^{-7t/10}+e^{-2t/5},
 \qquad B(t)=2e^{-t/2}+t e^{-3t/10}.
\]
This envelope bounds the sine term. Its derivative is
$B'(t)=-e^{-t/2}+(1-3t/10)e^{-3t/10}$.
For $t\geq10/3$ it is negative. For $0\leq t<10/3$,
$1-3t/10\leq1-t/5\leq e^{-t/5}$, so it is nonpositive.
Consequently $|H(t)|\leq B(0)+3/2+1=9/2$, attained at $t=0$.
In the exact-real interpretation the sample includes this maximum, hence
$H_{\max}=9/2+10^{-12}$ and
\begin{equation}\label{eq:gainbounds}
 \frac{99}{100}\leq a_i(t)\leq\frac{101}{100},
 \qquad \frac{19}{20}\leq b_t\leq\frac{21}{20}\quad(t\geq0).
\end{equation}
This is a reused proof about real formulas. It is not fresh floating-point
certification of the historical sampled normalizer.

\subsection{The common signed compressor}\label{sec:compression}
For $\epsilon=10^{-9}$ and common $L$, the amplitude compressor is
\begin{equation}\label{eq:compressor}
 C_L(q)=
 \begin{cases}
 0,&q=0,\\
 \sgn(q)\min\{e^L,\max\{e^{-L},|q|+\epsilon\}\},&q\ne0.
 \end{cases}
\end{equation}
Ordinarily $L=9$. With the optional common breathing enabled, the source instead uses
\begin{equation}\label{eq:breathing}
 L=9\left(1+\frac12\,\operatorname{mean}_i|R_i|\right),
\end{equation}
computed from the incoming amplitudes, once for the whole pass.
It does not create a different compressor for each particle.

For nonzero input of fixed sign, the floor fiber is
$0<|q|\leq e^{-L}-\epsilon$, provided this upper bound is positive.
It is empty if $e^{-L}\leq\epsilon$. The ceiling fiber is
$|q|\geq e^L-\epsilon$, and between these boundaries the map adds
$\epsilon\sgn(q)$. Endpoints belong to the plateaus.
In particular, a large breathing value must not be assumed to leave a positive-width
floor fiber. Zero is a separate branch and has output zero, not a signed floor.

\section{Equality is conditional, and compression matters}\label{sec:equality}
\begin{proposition}[Exact criterion]\label{prop:criterion}
Suppose $R_i=r$ for every particle. With the common $L$ chosen by the source,
\begin{equation}\label{eq:coeff}
 q_i=b_t r c_i,\qquad
 c_i=
 \begin{cases}
 a_i(t)+\sum_{j\in N_i}\eta_{ij},&i\text{ active},\\
 1,&i\text{ skipped}.
 \end{cases}
\end{equation}
The next amplitude vector is homogeneous if and only if all numbers
$C_L(b_t r c_i)$ agree.
\end{proposition}
\begin{proof}
On each active row, normalization gives $\sum_j w_{ij}r=r$.
Its mixed baseline is therefore $(1-\beta_T)r+\beta_T r=r$.
Equation~\eqref{eq:full} gives \eqref{eq:coeff}; the source's final amplitude
readout is exactly $C_L(q_i)$. This proves both directions of the criterion.
\end{proof}
Equal $c_i$ are sufficient, but not necessary: distinct inputs of the same sign
can occupy one floor or ceiling plateau. On a common unclipped branch, the compressor
is injective. This distinction prevents a raw-amplitude difference from being
mistaken for an output difference without a branch proof.

\begin{theorem}[Equal-frequency active rows]\label{thm:equal}
Let every incoming amplitude be $r$ and every frequency be the same positive $\omega$.
If every row is active, then
\begin{equation}\label{eq:equal}
 R_i^{\rm next}=C_L\!\left(b_t a_\omega(t)r\right)
 \quad\hbox{for every }i.
\end{equation}
This holds for arbitrary finite incoming distances and normalized weights for which
the processing hypotheses hold, including the common breathing choice.
\end{theorem}
\begin{proof}
For equal frequencies, $|\omega_i-\omega_j|=0$ closes the outer strict
enhancement guard, so every $\eta_{ij}=0$; the ratio terms do not survive that
guard. Also $a_i=a_\omega$. Proposition~\ref{prop:criterion} then gives a common
precompression value and a common output.
If $r=0$, all terms vanish and $C_L(0)=0$.
If $r>0$, the bounds \eqref{eq:gainbounds} make $b_ta_\omega r>0$ and the output
is positive; if $r<0$ both the precompression and final values are negative.
Thus neither a lost sign nor a zero-to-floor convention is hidden in the conclusion.
\end{proof}

\begin{proposition}[Conditional induction]\label{prop:induction}
Begin with homogeneous amplitudes and common positive frequencies.
For any finite sequence of passes on which the source remains valid and every row
is active, amplitudes remain homogeneous after each pass.
\end{proposition}
\begin{proof}
Apply Theorem~\ref{thm:equal} to the first pass. Its common output is the common
incoming amplitude for the second. The same theorem applies again because
frequency equality and active processing are assumed at every pass. Induction
proves the assertion, with the common breathing value recomputed if enabled.
\end{proof}
No assertion here proves that a moving particle system keeps all rows active forever.
The common scalar may vary with the clock. Equality of amplitudes is not constancy
of that scalar and says nothing about equality of positions or velocities.

\section{The minimal geometric active/skip split}\label{sec:geometric}
Return to the three particles in Section~\ref{sec:opening}. Fix $L=9$ with breathing
off, $r=1/5$, $\omega=61/250$, $\rho=1$ and $t=1$.
The distances are $1/2$, $5/2$ and $3$, so
\begin{equation}\label{eq:threegraph}
 N_1=\{2\},\qquad N_2=\{1\},\qquad N_3=\varnothing.
\end{equation}
The sole raw weight of each active row is
$1/(1/2+10^{-8})=100000000/50000001>10^{-10}$.
All enhancement vanishes by frequency equality. The active and skipped
precompression values are respectively
$q_{\rm pair}=b_1r a_\omega(1)$ and $q_{\rm iso}=b_1r$.

The eligibility, separation and activity statements are strict. Each position
can be perturbed by norm less than $1/4$: pairwise distances change by less
than $1/2$, the close pair stays distinct and within the horizon, and the other
distances remain outside it. The close-pair raw sum remains greater than
$1/(1+10^{-8})$. This is a local geometric margin, not a claim about a future trajectory.

Both precompression values satisfy
\begin{equation}\label{eq:threebounds}
 \frac{1881}{10000}\leq q\leq\frac{2121}{10000}.
\end{equation}
Indeed these are $r(19/20)(99/100)$ and $r(21/20)(101/100)$,
and also bound the skipped value. Since $e>2$,
$e^{-9}<1/512<1881/10000$, while
$2121/10000+10^{-9}<512<e^9$.
Thus both outputs are $q+\epsilon$, strictly inside the positive unclipped branch.
Their difference is exactly
\begin{equation}\label{eq:threegap}
 \Delta_{\rm geom}=b_1\,\frac15\,\frac1{100}
                   \sin\!\left(\frac{61\pi}{125}\right)>0.
\end{equation}
The sine argument lies in $(0,\pi)$ and $b_1>0$. In particular, $H$ multiplies
the isolated particle too; omitting that stage would give the wrong difference.

\begin{proposition}[Minimality for the common-frequency mixed-processing mechanism]
\label{prop:threeminimal}
One or two particles cannot realize a common-frequency active/skip split under
the common selection parameters of the source. Three suffice.
\end{proposition}
\begin{proof}
One particle provides no internal pair to compare. For two distinct particles,
their mutual distance and equal-frequency raw weight agree in both directions.
With a common horizon they are either both eligible or both ineligible.
If eligible, the same raw sum makes both active or both skipped. Without a horizon,
each selects the other and their processing agrees. At a coincident tie, a sorted
slice may instead select self; with common amplitudes and frequencies the same
zero distance, baseline and guard give no active/skip distinction.
Hence two cannot realize this mechanism. Equation~\eqref{eq:threegraph} supplies
the three-particle witness.
\end{proof}

The accepted corrected OC05 value is
\begin{equation}\label{eq:threenumeric}
 \Delta_{\rm geom}\approx0.002058402961068474
 \quad\hbox{(rounded to 18 decimal places).}
\end{equation}
The initially supplied shorter decimal failed a comparison and was corrected in
OC05. That initial failure and the accepted correction are preserved, not replaced
by a retrospective all-pass account; Section~\ref{sec:evidence} records the provenance.

\section{Two active particles with the actual two-band assignment}\label{sec:bands}
The second mechanism changes frequency, not eligibility. Specify an incoming state
of two particles at $0$ and $1/2$, with amplitudes $r=1/5$, no horizon,
no breathing, $L=9$ and incremented clock $1$. This is a prescribed incoming
state, not a claim that a randomized initializer samples it.

\begin{center}
\begin{tabular}{@{}ccccc@{}}
\toprule Role & Position & Source band index & Frequency & Neighbor\\
\midrule
low & $(0,0,0)$ & $-1$ & $\Omega/\phi$ & high\\
high & $(1/2,0,0)$ & $0$ & $\Omega$ & low\\
\bottomrule
\end{tabular}
\end{center}

The source truncates its centered index list to the requested band count.
For two bands the list is exactly $[-1,0]$, not a guessed symmetric pair:
\begin{equation}\label{eq:bandfreq}
 \phi=\frac{1+\sqrt5}{2},\qquad \Omega=\frac{61}{250},
 \qquad \omega_{\rm low}=\Omega/\phi,\quad\omega_{\rm high}=\Omega.
\end{equation}
Shuffling labels exchanges roles but does not change the argument.

\subsection{Strict activity and the complete enhancement}
Put $\ell=\log\phi$. Since $3/2<\phi<2$, we have $0<\ell<1$.
Each row has exactly one neighbor and
\begin{equation}\label{eq:bandactive}
 S_i=\frac{e^{-3\ell^2}}{1/2+10^{-8}}
        >e^{-3}>\frac1{27}>10^{-10},
\end{equation}
using $e<3$. Thus normalization yields the incoming amplitude $r$ in each row.
The frequency difference and sum are
\begin{equation}\label{eq:bandguard}
 D=\Omega(1-\phi^{-1})=\Omega\phi^{-2}
   >\frac{61}{750}>\frac1{100},
 \qquad U=\Omega(1+\phi^{-1})=\Omega\phi.
\end{equation}
The strict outer guard is open, with lower margin
$61/750-1/100=107/1500$.

The tower list has just the two frequencies in \eqref{eq:bandfreq}.
Every sum/difference indicator in Appendix~\ref{app:eta} fails:
\begin{center}
\begin{tabular}{@{}ccc@{}}
\toprule Tower entry $\tau$ & $|U-\tau|/\tau$ & $|D-\tau|/\tau$\\
\midrule
$\Omega/\phi$ & $\phi$ & $\phi^{-2}$\\
$\Omega$ & $\phi^{-1}$ & $\phi^{-1}$\\
\bottomrule
\end{tabular}
\end{center}
All four values exceed $1/5$: $\phi>3/2$, $\phi^{-1}>1/2$ and
$\phi^{-2}>1/4$. No equality-at-tolerance convention is being used.
For ratio powers $p=-2,-1,1,2$, the squared multiples of $\ell$ in the Gaussian
are $(1,0,4,9)$ for low receiving high and $(9,4,0,1)$ for high receiving low.
Therefore both complete coefficients equal
\begin{equation}\label{eq:etastar}
 \eta_*=\frac{1+e^{-5\ell^2}+e^{-20\ell^2}+e^{-45\ell^2}}{1000},
 \qquad \frac1{1000}<\eta_*<\frac4{1000}.
\end{equation}
Enhancement is present and positive. It has not been discarded to obtain a
sine-only model.

\subsection{A strict output difference with enhancement retained}
The precompression values are
\begin{equation}\label{eq:bandq}
 q_{\rm low}=b_1r(a_{\rm low}+\eta_*),\qquad
 q_{\rm high}=b_1r(a_{\rm high}+\eta_*).
\end{equation}
Equations~\eqref{eq:gainbounds} and \eqref{eq:etastar} give the common enclosure
\begin{equation}\label{eq:bandbranch}
 \frac{1881}{10000}<q_i<\frac{10647}{50000}.
\end{equation}
As in the geometric example, $e^{-9}<1/512$ is below the lower bound, and
$10647/50000+10^{-9}<512<e^9$ is above the shifted upper bound.
Both outputs are positive and unclipped.

Moreover $0<\omega_{\rm low}<\omega_{\rm high}<1/4$, with
$1/4-\Omega=3/500$. Sine is strictly increasing between the corresponding
arguments in $(0,\pi/2)$. Subtracting the two compressed values gives
\begin{equation}\label{eq:bandgap}
 \Delta_{\rm band}=\frac{b_1r}{100}
 \left[\sin(2\pi\Omega)-\sin(2\pi\Omega/\phi)\right]>0.
\end{equation}
The identical enhancement terms and identical positive-branch $\epsilon$ additions
cancel in this difference. Their cancellation is specific to the proven fixture.
In particular, $b_1r\eta_*>0$ raises both output levels.

\begin{center}
\begin{tabular}{@{}lr@{}}
\toprule Accepted OC07 quantity & Saved value, rounded to 18 decimal places\\
\midrule
$\eta_*$ & $0.001323940975761039$\\
$b_1r\eta_*$ & $0.000272714173360112$\\
$R_{\rm low}^{\rm next}$ & $0.207931904558321191$\\
$R_{\rm high}^{\rm next}$ & $0.208317773395512482$\\
$\Delta_{\rm band}$ & $0.000385868837191291$\\
\bottomrule
\end{tabular}
\end{center}

One particle cannot exhibit an inter-particle difference, while these two do;
this proves minimality for the established all-active frequency mechanism.
It does not contradict Proposition~\ref{prop:threeminimal}, whose frequencies
are equal and whose proposed mechanism is mixed processing.
With a sole neighbor per row, normalized averaging has weight one. Varying their
distance while both rows stay active does not change this calculation.
The example is not evidence for distance-induced memory recovery.

\section{What equality does and does not establish}\label{sec:limits}
Common positive amplitude need not remain at the positive floor. For example, take
the common value $r=e^{-9}$, fixed $L=9$, equal frequencies and all rows active.
With $A_t=b_ta_\omega(t)>0$, staying at that floor is equivalent to
\begin{equation}\label{eq:floorstay}
 A_t e^{-9}+\epsilon\leq e^{-9}
 \quad\Longleftrightarrow\quad A_t\leq1-\epsilon e^9 .
\end{equation}
Failure of this inequality may raise every amplitude together.
Conversely, plateaus can make different precompression values agree.
These are different statements from full-state capture or divergence.

The spatial update, fusion and positional noise remain part of the historical
composition. They can alter later geometry without altering a same-pass amplitude
readout already determined by stored incoming neighborhoods. Nothing proved here
identifies positions, velocities or all later neighbor choices.
The specified fixtures are not later passes of the accepted OC04 geometry witness. The whole-state qualification is also retained in Geometry~\cite{Geometry}.

The accepted scope comprises exact-real one-pass identities, conditional induction,
strict local branch proofs and minimal examples for two stated mechanisms.
It does not establish stochastic arrival, autonomous heartbeat, entropy production,
cosmology, physical validation or general long-time behavior.
The separate January fixed-grid question concerns access to signed floor targets;
its states, gains and constants are not premises of these Tower proofs.

\section{Evidence, authorship and review status}\label{sec:evidence}
OC05 and OC07 are accepted within their declared analysis scopes. This new v0.1
manuscript and its six-file proposed public payload await manuscript-and-payload
review. No external peer review or publication acceptance is asserted.
The human author is Hilmir Frímann Halldórsson. GPT supplied the work orders,
scientific review and candidate deductions in the accepted analyses; Codex prepared
this exposition, performed static source and proof correspondence checks, transcribed
saved decimals, and built and inspected the manuscript. These AI activities are
assistance, not independent human peer review.

The decimal in \eqref{eq:threenumeric} comes from OC05's accepted corrected saved value:
\[
 0.002058402961068473853043\ldots .
\]
The initial comparison against $0.002058402961068466$ failed.
The retained initial record has five passing groups and one failed group.

The final accepted OC05 record has nine passing groups.
OC07 has eight passing groups. Those are historical result counts, not checks
rerun for this manuscript. No accepted checker, supplied review script, historical
program or rollout was executed during preparation.

The rounded examples use 18 decimal places with decimal round-half-even.
The correction-history digits above are retained as quoted provenance; no transcendental
value was reevaluated.
The exact formulas and inequalities carry the proofs. The accompanying source index
and claim ledger record versions, physical line references and accepted-report
attribution, including private-source access limitations.

The main conclusion is conditional and concrete: common active processing with
equal frequencies preserves amplitude equality; an isolated row can remove the sine
stage, and different frequencies can split active rows even when their enhancement
increments agree. Compression must be checked in either example.
No new blanket reuse licence is assigned to this manuscript or reconstruction
records. Existing source-specific CC BY 4.0 notices remain specific to those original
materials; a software licence is not presumed to license this research.

\appendix
\section{The complete source enhancement coefficient}\label{app:eta}
Let $\mathcal T$ be the actual tower-frequency list, with multiplicities retained,
and set
$U_{ij}=\omega_i+\omega_j$, $D_{ij}=|\omega_i-\omega_j|$.
For a selected directed pair the source coefficient is
\begin{equation}\label{eq:eta}
\begin{split}
 \eta_{ij}={}&\nu\,\ind_{\{D_{ij}>0.01\}}
 \Bigg[
 \sum_{p\in\{-2,-1,1,2\}}
 \exp\!\left(-5[\log(\omega_i/\omega_j)-p\log\phi]^2\right)\\
 &\quad+\frac12\sum_{\tau\in\mathcal T}
 \left(
 \ind_{\{|U_{ij}-\tau|/\tau<0.2\}}
 +\ind_{\{D_{ij}>0.01\}}
  \ind_{\{|D_{ij}-\tau|/\tau<0.2\}}
 \right)
 \Bigg],\qquad \nu=0.001.
\end{split}
\end{equation}
The repeated difference guard reflects the nested source test, even though the
outer guard already implies it. Each successful tower-list test contributes;
the list is not a set of distinct values. Every tolerance is strict.
The coefficient multiplies the incoming neighbor amplitude $R_j$ and is summed
after the sine stage. No distance denominator or normalization by $S_i$ belongs
in \eqref{eq:eta}. The shared coefficient in \eqref{eq:etastar} was obtained by
checking all these terms for the actual two-entry tower.

\section{Source line map and references}\label{app:sources}
Line numbers are one-based in the frozen, complete Tower text T; byte identity
and archive locators are supplied in \texttt{SOURCE\_INDEX.md}.
\begin{center}\small
\begin{tabularx}{\textwidth}{@{}lY@{}}
\toprule T lines & Definition used\\
\midrule
31--46 & $\phi$, $\Omega$, $H$, RPCO and sampled $H$ normalization\\
62--84 & defaults, actual band-index construction and shuffled labels\\
93--103 & $L$, spatial and amplitude coefficients, $H$ coefficient, neighbor count\\
111--140 & selector, clock increment, common breathing, incoming distances\\
142--167 & spatial motion, row guards, normalized mixing and sine\\
169--194 & complete enhancement followed by common $H$\\
197--215 & precompression positional fusion, RPCO, noise and velocity\\
\bottomrule
\end{tabularx}
\end{center}
The accepted source-order addendum A, lines 33--47, qualifies the skip and order
statements. G, lines 185--214, supplies the inherited $H$ bound. Central arguments
are reproduced here so that private-report access is not required to follow them.

\begin{thebibliography}{9}
\bibitem{T} Historical source, \emph{srg\_golden\_tower\_II.py},
14 March 2026. Frozen text T, 34,980 bytes. Header credits Claude collaborating
with pzychozen. Original public archive locator:
\url{https://zenodo.org/records/19023429}, \texttt{simulations.zip}.
Exact member and SHA-256 are in the companion source index.
\bibitem{OC05} \emph{OC05 Tower State Persistence}, accepted REPORT.md and
CHECK\_RESULTS.json, 12 October 2026. Retained private research records,
including initial failed-decimal comparison and accepted correction.
\bibitem{OC07} \emph{OC07 Tower Two-Band Split}, accepted REPORT.md and
CHECK\_RESULTS.json, 12 October 2026. Retained private research records.
\bibitem{G} \emph{OC02 Signed Plateaus}, accepted report G, 11 October 2026.
Retained final handoff, \path{updates/OC02_SIGNED_PLATEAUS_REPORT.md}.
\bibitem{A} \emph{Required Source and Notation Addendum}, retained final handoff,
\path{updates/REQUIRED_SOURCE_AND_NOTATION_ADDENDUM.md}.
\bibitem{Geometry} Hilmir Frímann Halldórsson, \emph{SRG Geometry: Local Response, Neighbor Exchange, and Compression}, v0.1,
12 October 2026, published local edition, Section 6.2, background on floor and
whole-state qualifications. Source access routes are inherited from that edition;
no new remote query was made.
\end{thebibliography}
\end{document}
